\documentclass[10pt,a4paper]{article}
\usepackage[margin=1.4cm, top=1.2cm, bottom=1.2cm]{geometry}
\usepackage{booktabs}
\usepackage{array}
\usepackage{xcolor}
\usepackage{colortbl}
\usepackage{caption}
\usepackage[skip=3pt]{parskip}
\usepackage{titlesec}
\usepackage{microtype}
\usepackage{amsmath}

\definecolor{headerblue}{RGB}{30,90,160}
\definecolor{rowgray}{RGB}{240,240,240}

\titleformat{\section}{\normalfont\large\bfseries\color{headerblue}}{}{0em}{}
\titlespacing{\section}{0pt}{4pt}{2pt}

\captionsetup{font=small, labelfont=bf, skip=4pt}

\begin{document}

\begin{center}
{\Large \textbf{Assignment 2 --- Problem 1: Loop Re-ordering (Cache Effects)}}\\[3pt]
{\normalsize CS60003: High Performance Computer Architecture \quad|\quad IIT Kharagpur \quad|\quad Spring 2025--26}
\end{center}

\vspace{2pt}
\hrule
\vspace{4pt}

\section{Experimental Setup}

An $N \times N$ matrix of \texttt{double}-precision floats (declared globally) was filled with
random values. Two traversal orders were implemented: \textbf{row-major} (inner loop over columns,
cache-friendly in C) and \textbf{column-major} (inner loop over rows, strided access). Each was
compiled with \texttt{-O0} (no optimisation) and \texttt{-O3 -march=native -fopt-info-vec}
(aggressive optimisation + vectorisation). All experiments were pinned to \textbf{CPU core 2}
(P-core) using \texttt{taskset} to eliminate scheduling noise. Cache statistics were collected
using \texttt{perf stat} on an Intel hybrid CPU; only \texttt{cpu\_core} counters are reported
as \texttt{cpu\_atom} counters were \texttt{<not counted>} on P-core~2.

\vspace{4pt}
\section{Results}

\begin{table}[h!]
\centering
\small
\setlength{\tabcolsep}{6pt}
\renewcommand{\arraystretch}{1.15}

\caption{Run time, cache accesses, cache misses, and miss rate for all 16 configurations (pinned to P-core 2).}

\begin{tabular}{c l l r r r}

\toprule
\rowcolor{headerblue}
\color{white}\textbf{N} &
\color{white}\textbf{Access Type} &
\color{white}\textbf{Opt. Flag} &
\color{white}\textbf{Time (s)} &
\color{white}\textbf{Cache References} &
\color{white}\textbf{Cache Misses (Miss \%)} \\
\midrule

1024 & Row-major    & \texttt{-O0} & 0.002077 & 316,419   & 186,612 (58.98\%) \\
     & Row-major    & \texttt{-O3} & 0.000624 & 315,324   & 221,117 (70.12\%) \\
\rowcolor{rowgray}
     & Column-major & \texttt{-O0} & 0.003069 & 1,233,402 & 233,219 (18.91\%) \\
\rowcolor{rowgray}
     & Column-major & \texttt{-O3} & 0.003381 & 1,235,328 & 280,484 (22.71\%) \\

\midrule

2048 & Row-major    & \texttt{-O0} & 0.008004 & 1,200,533  & 883,349   (73.58\%) \\
     & Row-major    & \texttt{-O3} & 0.002618 & 1,183,563  & 1,049,764 (88.70\%) \\
\rowcolor{rowgray}
     & Column-major & \texttt{-O0} & 0.021900 & 4,862,554  & 1,744,482 (35.88\%) \\
\rowcolor{rowgray}
     & Column-major & \texttt{-O3} & 0.020642 & 4,882,134  & 1,587,729 (32.52\%) \\

\midrule

4096 & Row-major    & \texttt{-O0} & 0.031492 & 4,678,029  & 3,423,735  (73.19\%) \\
     & Row-major    & \texttt{-O3} & 0.009083 & 4,670,829  & 4,135,685  (88.54\%) \\
\rowcolor{rowgray}
     & Column-major & \texttt{-O0} & 0.164405 & 19,640,191 & 15,367,724 (78.25\%) \\
\rowcolor{rowgray}
     & Column-major & \texttt{-O3} & 0.164094 & 19,606,503 & 15,088,652 (76.96\%) \\

\midrule

8192 & Row-major    & \texttt{-O0} & 0.126813 & 18,293,841 & 13,400,024 (73.25\%) \\
     & Row-major    & \texttt{-O3} & 0.036263 & 18,141,914 & 16,484,855 (90.87\%) \\
\rowcolor{rowgray}
     & Column-major & \texttt{-O0} & 0.770434 & 88,231,398 & 70,113,874 (79.47\%) \\
\rowcolor{rowgray}
     & Column-major & \texttt{-O3} & 0.748653 & 86,655,856 & 69,297,092 (79.97\%) \\

\bottomrule
\end{tabular}

\end{table}

\vspace{1pt}

\section{Observations}

\textbf{1. Row-major is consistently faster than column-major.}
At $N=8192$, row-major (\texttt{-O0}: 0.127\,s, \texttt{-O3}: 0.036\,s) is $6.07\times$ and
$20.65\times$ faster than column-major (0.770\,s and 0.749\,s) respectively. C stores 2-D arrays
in row-major order, so consecutive row elements are contiguous in memory. Row-major traversal
exploits this spatial locality, loading a full 64-byte cache line per step. Column-major traversal
strides by $N \times 8$ bytes between accesses, causing a cache miss on nearly every element.

\textbf{2. Speedup gap widens with increasing N.}
At $N=1024$ row-major is $\sim1.48\times$ faster (\texttt{-O0}), rising to $\sim6.07\times$ at
$N=8192$. As $N$ grows the matrix no longer fits in L2/L3 cache. Row-major still benefits from
hardware prefetching on sequential addresses, while column-major suffers increasingly severe
cache-line evictions between consecutive accesses in the inner loop.

\textbf{3. \texttt{-O3} benefits row-major far more than column-major.}
The compiler auto-vectorises the row-major inner loop (confirmed by \texttt{loop vectorized using
32 byte vectors} in compiler output), issuing 256-bit AVX2 SIMD instructions that process 4
\texttt{double} values per cycle. Column-major access cannot be vectorised effectively due to its
non-contiguous memory pattern, so \texttt{-O3} yields minimal speedup for column-major
(e.g., $0.164\,\text{s} \to 0.164\,\text{s}$ at $N=4096$, and only $0.770\,\text{s} \to
0.749\,\text{s}$ at $N=8192$).

\textbf{4. Row-major generates far fewer cache references than column-major.}
At $N=4096$ with \texttt{-O0}, row-major generates only 4.7M references vs.\ column-major's
19.6M --- a $4\times$ difference. Row-major's sequential pattern allows the hardware prefetcher
to silently satisfy many accesses before they stall, so they never appear as cache reference
events. Column-major defeats the prefetcher entirely, generating far more visible cache traffic
and more total misses in absolute terms, directly explaining its slower execution time.

\textbf{5. Miss rate anomaly at small N and its explanation.}
At $N=1024$, row-major shows a higher miss \emph{rate} (58.98\%) than column-major (18.91\%).
This is not a contradiction: row-major's reference count is $4\times$ smaller because the
prefetcher hides most accesses from the counter. The \emph{absolute} miss count is still lower
for row-major (186K vs.\ 233K), and execution time confirms row-major is faster (0.002\,s
vs.\ 0.003\,s). At $N \geq 4096$ both patterns show high miss rates ($>73\%$) as the working
set exceeds L3 cache capacity, but column-major consistently incurs more total misses and
$5\text{--}20\times$ slower runtime.

\end{document}
